\section{Objek: Literal, Properti, dan Metode}

Objek adalah kumpulan pasangan \textbf{key-value} yang merepresentasikan entitas dengan properti dan perilaku. Objek dibuat dengan literal \code{\{\}} dan diakses dengan dot notation (\code{obj.key}) atau bracket notation (\code{obj["key"]}). Properti menyimpan data; metode adalah fungsi yang menjadi properti objek. Objek adalah tipe data paling fleksibel di JavaScript---hampir semua hal di JS adalah objek \cite{jsInfo}.

\begin{webcode}[language=JavaScript, caption={Objek literal, properti, dan metode}]
const mahasiswa = {
  nama: "Budi Santoso",
  nim: "2024001",
  nilai: [85, 90, 78],
  // Metode (fungsi di dalam objek)
  rataRata() {
    const total = this.nilai.reduce((s, n) => s + n, 0);
    return total / this.nilai.length;
  }
};

console.log(mahasiswa.nama);      // "Budi Santoso"
console.log(mahasiswa["nim"]);    // "2024001"
console.log(mahasiswa.rataRata()); // 84.33...

// Destructuring objek
const { nama, nim } = mahasiswa;
console.log(nama); // "Budi Santoso"
\end{webcode}

\begin{table}[htbp]
\centering
\begin{tabular}{|P{3.5cm}|P{4cm}|P{4.5cm}|}
\hline
\textbf{Operasi} & \textbf{Sintaks} & \textbf{Keterangan} \\
\hline
Akses properti & \code{obj.key} / \code{obj["key"]} & Bracket untuk key dinamis \\
\hline
Tambah properti & \code{obj.newKey = val} & Dinamis; kapan saja \\
\hline
Hapus properti & \code{delete obj.key} & Menghapus dari objek \\
\hline
Cek keberadaan & \code{"key" in obj} & Mengembalikan boolean \\
\hline
Iterasi keys & \code{Object.keys(obj)} & Mengembalikan array of keys \\
\hline
Spread & \code{\{...obj, newKey: val\}} & Copy + tambah properti (ES6) \\
\hline
\end{tabular}
\caption{Operasi umum pada objek JavaScript}
\end{table}

\begin{catatan}
\textbf{Object vs Array.} Gunakan \textbf{array} untuk koleksi terurut dari item serupa (daftar nama, nilai ujian). Gunakan \textbf{objek} untuk entitas dengan properti bernama yang berbeda (mahasiswa, produk, konfigurasi). Array memiliki indeks numerik dan metode iterasi; objek memiliki key string/symbol dan cocok untuk lookup \cite{mdnJsGuide}.
\end{catatan}

